\documentclass[10pt,a4paper]{amsart}
\usepackage[margin=19mm]{geometry}
\usepackage{amsmath,amssymb,mathtools}
\usepackage[scr=rsfs,cal=euler]{mathalfa}
\usepackage{xfrac}
\usepackage[unicode,hidelinks,bookmarks=false]{hyperref}
\newcommand{\ie}{i.e.\ }
\newcommand{\cf}{cf.\ }
\newcommand{\resp}{respectively\ }
\newcommand{\Subsection}{\subsection}
\providecommand{\nobreakdash}{\nobreak\hbox{-}}
\setlength{\emergencystretch}{2em}
\multlinegap0pt
\newcommand{\I}{\mathbf{1}}
\newcommand{\nuc}[1]{\left\lVert #1 \right\rVert_{*}}
\theoremstyle{plain}
\newtheorem{theorem}{Theorem}
\newtheorem{proposition}[theorem]{Proposition}
\newtheorem{corollary}[theorem]{Corollary}
\theoremstyle{remark}
\newtheorem{remark}[theorem]{Remark}
\title[A sharp norm inequality for entanglement-breaking channels]{A sharp norm inequality for entanglement-breaking channels}
\author{Eran Kopel}
\date{Source: arXiv:2609.27906v1 (CC BY 4.0)}
\begin{document}
\maketitle

\noindent\emph{Source and licence.} The delivered work is the article shown in the title above, version arXiv:2609.27906v1 (CC BY 4.0), distributed under the Creative Commons Attribution 4.0 International licence, as recorded in the arXiv licence metadata for this exact version and shown on the abstract page saved with this item. The full licence notice, which carries the licence URL and the address of that abstract page, is the bundled file \texttt{licenses/arxiv-2609.27906-CC-BY-4.0.txt}.
\par\smallskip

\noindent\textit{Editorial scope.} Counted unit: the article's Theorem (source label thm:main) (source0.tex lines 82--88, source label \texttt{thm:main}): ``Let be an entanglement-breaking channel on with Bloch data . Then A 2 ; ; d(d-1) 2 , c 2 ; ; (d-1) 2 . For this reads .''. It is delivered together with the article's complete original proof (source0.tex lines 90--124, one proof environment adjacent to the statement, 35 lines), copied line for line from the article's own text. Required context retained uncounted: eq:bloch (lines 69--71); eq:povm (lines 73--77). Editorial formatting only: an AMS \texttt{amsart} preamble that restates the article's own macro definitions and its own theorem declarations, and the theorem environments the retained lines use; the article's own class file is neither shipped nor input; the retained environments are renumbered sequentially in order of appearance, whereas the article prints them with its own numbering; the article's equation numbering is not reproduced and every cross-reference inside the retained text resolves to the wrapper's numbering. No citation occurs inside any retained line, so the wrapper carries no bibliography; All retained bodies are exact copies of the bundled \texttt{sources/211525/source0.tex}, so no omission receipt applies. No asset-inclusion command occurs inside any retained span and the package ships no image asset.


\section*{Retained context: eq:bloch (source0.tex lines 69--71)}

\begin{equation}\label{eq:bloch}
A \;=\; \tfrac12\sum_k s_k f_k^{\mathsf T}, \qquad c \;=\; \sum_k \tfrac{\alpha_k}{d}\,s_k .
\end{equation}


\section*{Retained context: eq:povm (source0.tex lines 73--77)}

\begin{equation}\label{eq:povm}
\sum_k E_k=\I \;\Longrightarrow\; \sum_k\alpha_k=d \ \text{ and }\ \sum_k f_k=0,
\qquad
E_k\succeq0 \;\Longrightarrow\; \Big|\tfrac{f_k}{\alpha_k}\Big|^2 \le \tfrac{2(d-1)}{d},
\end{equation}


\section{The counted unit: Theorem (source label thm:main) and its complete original proof}

\begin{theorem}\label{thm:main}
Let $\Phi$ be an entanglement-breaking channel on $M_d$ with Bloch data $(A,c)$. Then
\[
\nuc{A}^2 \;+\; \frac{d(d-1)}{2}\,|c|^2 \;\le\; (d-1)^2 .
\]
For $d=2$ this reads $\nuc{A}^2+|c|^2\le1$.
\end{theorem}

\begin{proof}
Put $u_k := \alpha_k/d \ge 0$ and $g_k := \tfrac{d}{2}\,f_k/\alpha_k$. By \eqref{eq:povm}, $\sum_k u_k=1$, so $\{u_k\}$ is a probability distribution, and
\begin{equation}\label{eq:gbound}
|g_k| \;\le\; \frac{d}{2}\sqrt{\frac{2(d-1)}{d}} \;=\; \sqrt{\frac{d(d-1)}{2}} \;=:\; G_d .
\end{equation}
In this notation \eqref{eq:bloch} becomes $A=\sum_k u_k s_k g_k^{\mathsf T}$ and $c=\sum_k u_k s_k$, while completeness reads
\begin{equation}\label{eq:centred}
\sum_k f_k=0 \quad\Longleftrightarrow\quad \sum_k u_k g_k = 0 .
\end{equation}

\emph{Step 1 (recentring).} By \eqref{eq:centred}, $\sum_k u_k\,c\,g_k^{\mathsf T} = c\big(\sum_k u_k g_k\big)^{\mathsf T} = 0$, hence
\begin{equation}\label{eq:recentre}
A \;=\; \sum_k u_k\,(s_k-c)\,g_k^{\mathsf T}.
\end{equation}
This is the only use of completeness. It replaces the second moment $\mathbb{E}[s\,g^{\mathsf T}]$ by the cross-covariance $\mathbb{E}[(s-\mathbb{E}s)\,g^{\mathsf T}]$, and is what makes the constant attainable.

\emph{Step 2 (rank-one triangle inequality).} Since $\nuc{vw^{\mathsf T}}=|v|\,|w|$, applying the triangle inequality to \eqref{eq:recentre} and then \eqref{eq:gbound},
\[
\nuc{A} \;\le\; \sum_k u_k\,|s_k-c|\,|g_k| \;\le\; G_d\sum_k u_k\,|s_k-c| .
\]

\emph{Step 3 (Cauchy--Schwarz).} With respect to the weights $u_k$,
\[
\sum_k u_k|s_k-c| \;\le\; \Big(\sum_k u_k|s_k-c|^2\Big)^{1/2}.
\]

\emph{Step 4 (variance identity).} Since $c=\sum_k u_k s_k$,
\[
\sum_k u_k|s_k-c|^2 \;=\; \sum_k u_k|s_k|^2 - |c|^2 \;\le\; \frac{2(d-1)}{d}-|c|^2 .
\]
Combining,
\[
\nuc{A}^2 \;\le\; G_d^2\Big(\frac{2(d-1)}{d}-|c|^2\Big) \;=\; (d-1)^2-\frac{d(d-1)}{2}|c|^2 . \qedhere
\]
\end{proof}

\end{document}
